% --------------------------------------------------
% 2. Theoretische Grundlagen
% --------------------------------------------------
\section{Theoretische Grundlagen}

Dieses Kapitel legt das theoretische Fundament für die Fallstudie. Es beleuchtet zunächst den aktuellen Stand der Forschung, definiert die Rolle von \gls{e2etests} in der Qualitätssicherung, stellt anschließend die beiden zu vergleichenden Frameworks vor und erläutert abschließend die Methodik der Nutzwertanalyse.

\subsection{Stand der Forschung}
Die Evaluierung und der Vergleich von Testautomatisierungswerkzeugen sind Gegenstand zahlreicher Publikationen der angewandten Informatik. Während grundlegende Werke wie Liggesmeyer \cite{Liggesmeyer2009} oder Spillner und Linz \cite{spillner2024basiswissen} die theoretischen Konzepte des Softwaretests und der Testautomatisierung umfassend behandeln, fokussieren sich neuere Publikationen häufig auf den Architekturvergleich spezifischer Frameworks. Der Kontrast zwischen traditionellen \gls{webdriver}-basierten Ansätzen wie Selenium und modernen, browserinternen Werkzeugen wie Cypress wird in der Fachliteratur oft anhand von Metriken wie Ausführungsgeschwindigkeit, Stabilität (Flakiness) und \gls{api}-Design diskutiert. 

Eine Lücke besteht jedoch oft in der detaillierten Betrachtung dieser Tools innerhalb stark reglementierter Unternehmensnetzwerke und spezifischer Legacy-Infrastrukturen (wie z.\,B. ältere \gls{jdk}-Versionen oder strikte Proxy-Vorgaben). Die vorliegende Fallstudie knüpft an den bestehenden Forschungsstand an, erweitert ihn jedoch um eine systematische Bewertung, die exakt diese praxisnahen infrastrukturellen und organisatorischen Rahmenbedingungen der JUMO GmbH \& Co. KG in den Mittelpunkt stellt.

\subsection{Softwaretest und \gls{e2etest}ing}
Die Testautomatisierung ist ein fest etablierter Bestandteil moderner Softwarequalitätssicherung \cite{Liggesmeyer2009}. Während Unit- und Integrationstests isolierte Code-Abschnitte oder Schnittstellen auf technischer Ebene prüfen, testen \gls{e2etests} (\gls{e2e}-Tests) das gesamte System über alle Schichten hinweg -- von der Benutzeroberfläche bis zur Datenbank. Sie dienen insbesondere dazu, regressionskritische Nutzungsszenarien aus Sicht der Endanwender in einer realitätsnahen Umgebung zu überprüfen und dadurch zentrale Geschäftsprozesse abzusichern \cite{spillner2024basiswissen}. Da \gls{e2e}-Tests in der Regel komplexer in der Einrichtung und laufzeitintensiver sind, ist die Wahl eines robusten, wartbaren und in die Systemlandschaft passenden Automatisierungs-Frameworks entscheidend für die langfristige Akzeptanz im Entwicklungsalltag.

\subsection{Betrachtete Test-Frameworks: Selenium und Cypress}
Für den praktischen Vergleich dieser Arbeit stehen die Werkzeuge Selenium und Cypress im Mittelpunkt.

\textbf{Selenium} ist ein langjährig etabliertes Standardwerkzeug der browserbasierten Testautomatisierung. Es steuert den Browser über das native \gls{webdriver}-Protokoll an und bietet breite Integrationsmöglichkeiten in unterschiedliche Programmiersprachen, wie beispielsweise Java \cite{Selenium}. Diese hohe Flexibilität und Rückwärtskompatibilität erfordern jedoch im Gegenzug oftmals einen höheren manuellen Konfigurationsaufwand, insbesondere bei der Handhabung von asynchronen Ladezeiten und der Synchronisation des Testablaufs.

\textbf{Cypress} wird demgegenüber häufig als moderne, entwicklerzentrierte Alternative für Webanwendungen herangezogen. Im Gegensatz zu Selenium läuft Cypress direkt im Browser und führt Tests in derselben Event-Loop wie die Anwendung selbst aus \cite{Cypress}. Dies ermöglicht eine direkte Rückmeldung während der Testausführung, ein automatisches Warten auf \gls{dom}-Elemente (\textit{\gls{autowaiting}}) sowie sehr gute Debugging-Möglichkeiten. Da Cypress jedoch primär auf JavaScript beziehungsweise TypeScript ausgelegt ist, bedarf es bei der Integration in bestehende Java-Ökosysteme oder stark reglementierte Firmennetzwerke einer gesonderten Betrachtung.

Der Mehrwert der vorliegenden Arbeit liegt demnach nicht in einer rein allgemeinen Toolbeschreibung, sondern in der kontextbezogenen Bewertung dieser Frameworks für konkrete Webanwendungen unter realen infrastrukturellen Restriktionen im Unternehmensumfeld.

\subsection{Nutzwertanalyse als Bewertungsmethode}
Zur strukturierten und objektiven Entscheidungsfindung wird die Nutzwertanalyse verwendet. Diese Methode erlaubt es, mehrere Handlungsalternativen anhand systematisch gewichteter qualitativer und quantitativer Kriterien nachvollziehbar zu vergleichen \cite{Kuehnapfel2021}. Das Verfahren gliedert sich in vier zentrale Schritte: die Definition der relevanten Bewertungskriterien, deren Gewichtung auf Basis der betrieblichen Anforderungen, die Ermittlung der Zielerfüllung (Punktevergabe) durch prototypische Messungen sowie die finale Berechnung des Gesamtnutzwertes. Damit wird ein bestehender praktischer Bedarf mit einer methodisch fundierten Bewertungslogik verbunden, um eine kontextbezogene Entscheidungshilfe abzuleiten.